\subsection{Exact Finite-Angle Collision Foundations}

The marker modules begin with a two-body collision. Relative motion sets
the encounter rate, and conservation of momentum and energy determines the
outgoing velocities. This construction provides the basis for the collision
samplers and for the diffusion approximation introduced later.  Two incoming objects have
velocities \(\bm v_1,\bm v_2\), masses \(m_1,m_2>0\), relative velocity
\(\bm u=\bm v_1-\bm v_2\), and relative speed \(u=|\bm u|\).  A differential
cross section \(\dd\sigma/\dd\Omega\) gives the effective encounter area per
unit outgoing solid angle \(\dd\Omega\).  The local event rate is proportional
to
\[
  n_1 n_2\,u\,\frac{\dd\sigma}{\dd\Omega},
\]
where \(n_1,n_2\) are local number densities.  This factor explains the three
ingredients seen later in every exact sampler: find local partners, weight the
trial by relative speed and cross section, and draw an outgoing direction.

For an elastic event, the centre-of-mass velocity
\[
 \bm V=\frac{m_1\bm v_1+m_2\bm v_2}{m_1+m_2}
\]
is unchanged.  The outgoing relative velocity has the same length \(u\) but a
new direction \(\widehat{\bm n}\).  The unique outgoing velocities are
\[
 \bm v'_1=\bm V+\frac{m_2}{m_1+m_2}u\widehat{\bm n},
 \qquad
 \bm v'_2=\bm V-\frac{m_1}{m_1+m_2}u\widehat{\bm n}.
\]
These formulas conserve total momentum and kinetic energy for unequal as well
as equal masses.  Spherical variables are reconstructed only after this
three-dimensional collision has been completed.

\paragraph{Material SIDM.}
In a material self-interacting dark matter event, each marker represents
a sample of the continuous material distribution. The material law specifies
the cross section, often quoted per unit mass. Elastic scattering changes physical particle velocities while preserving
their rest masses. Each module specifies how numerical marker weights and
any unequal-weight remainder represent this event. Constant, power-law, and
Yukawa--Born laws have different dependences on \(u\) and scattering angle.
Their parameters therefore need the normalization and units of the selected
law.

\paragraph{Subhalo CDM encounters.}
In the subhalo operator, each marker represents one resolved extended
object.  The encounter area is derived from a softened Plummer gravitational
scattering model rather than a material cross section per mass.  An encounter
may scatter two complete subhaloes or, when the calibrated capture criterion
is met, replace them by one complete remnant.  The remnant has mass
\(m_1+m_2\), and the lost relative orbital energy is recorded as unresolved
internal energy.  This inelastic capture channel is not part of material
SIDM or the present SphDSMC module.

\paragraph{What the word exact means here.}
The pair kinematics above conserve the stated two-body invariants exactly.
The Monte Carlo sampler draws trial events from an upper rate, called a
majorant, and accepts each trial with the ratio of the physical rate to that
upper rate. A valid majorant makes this an unbiased sampling of the
represented discrete event process.  Numerical error still enters
through marker sampling, finite collision shells or cells, force resolution,
event cadence, rate bounds, finite run time, and random realization.  The
dense Eulerian operator described later uses the same word differently: it
enumerates every retained velocity-cell pair and configured angular node.
That label does not remove its angular quadrature, velocity boundary, remap,
or timestep error.

The Fokker--Planck operator appears only after these finite-angle constructions.
It represents the accumulated effect of many small deflections through
drift and diffusion coefficients. Its validity depends on the small-angle
regime, and its numerical accuracy requires convergence of those coefficients
and of the velocity space update.
